\subsection{零块半尺度、例外三扇区与 prime-ledger 的不可压缩裂口}

正文与附录已经分别给出：有限秩加法/异常账本不能在严格同态语义下忠实承载乘法骨架、六倍窗零点集合在绝对规模与相对密度上的半尺度律、以及 \(m=10\) 的 bin-fold 退化谱三扇区骨架 \((55,52,37)\) 与例外塔维数闭合。把这些接口重新并读后，可得到如下更紧的派生封装。

\begin{theorem}[有限异常账本不能忠实同态化乘法骨架]\label{thm:derived-finite-anomaly-ledger-no-faithful-multiplicative-skeleton}
\leanverified{paper\_derived\_finite\_anomaly\_ledger\_no\_faithful\_multiplicative\_skeleton}
设 \(\Lambda\) 为任意有限生成交换群。
不存在单射
\[
\Phi:\NN_{\ge 1}^{\times}\longrightarrow \ZZ\oplus \Lambda,
\qquad
\Phi(n)=(e(n),s(n)),
\]
同时满足
\[
\Phi(ab)=\Phi(a)+\Phi(b)
\qquad
(\forall a,b\ge 1),
\]
并且其第一坐标 \(e:\NN_{\ge 1}^{\times}\to \ZZ\) 亦为单射。

等价地，任何有限秩异常账本都不能在严格同态语义下忠实承载自然数乘法骨架。
\end{theorem}

\begin{proof}
反设这样的 \(\Phi\) 存在。由于加法在直积上逐坐标进行，\(\Phi(ab)=\Phi(a)+\Phi(b)\) 蕴含
\[
e(ab)=e(a)+e(b)
\qquad
(\forall a,b\ge 1),
\]
故 \(e\) 本身就是从乘法幺半群 \(\NN_{\ge 1}^{\times}\) 到加法群 \(\ZZ\) 的幺半群同态。
特别地，
\[
e(1)=0.
\]
由 \(e\) 单射可知，对任意素数 \(p\) 都有 \(e(p)\neq 0\)。

现取两个不同素数 \(p\neq q\)，并记
\[
a:=e(p),\qquad b:=e(q).
\]
若 \(ab>0\)，则
\[
e\!\bigl(p^{|b|}\bigr)=|b|a=|a|b=e\!\bigl(q^{|a|}\bigr),
\]
与 \(e\) 的单射性矛盾。
若 \(ab<0\)，则
\[
e\!\bigl(p^{|b|}q^{|a|}\bigr)=|b|a+|a|b=0=e(1),
\]
再度与单射性矛盾。
故这样的 \(e\) 不存在，从而这样的 \(\Phi\) 也不存在。

这与定理 \ref{thm:derived-cofinal-prime-support-unbounded-ledger-rank} 和定理
\ref{thm:derived-strict-multiplicative-ledger-binary-exhaustion}
所给出的“有限秩严格账本障碍”一致；这里只是把禁闭 sharpen 到了 \(\ZZ\oplus \Lambda\) 这一 rank-one 加有限 torsion 的最小异常账本模型。证毕。
\end{proof}

\begin{theorem}[六倍窗零块的半尺度双相律]\label{thm:derived-window6-zero-block-dual-halfscale-law}
\leanverified{paper\_derived\_window6\_zero\_block\_dual\_halfscale\_law}
设
\[
m+2=2d\equiv 0\pmod 6,
\qquad
M_m:=F_{m+2},
\qquad
\mathcal Z_m:=\{r:\widehat d_m(r)=0\}.
\]
则同时成立：
\[
\boxed{\;
|\mathcal Z_m|\ge F_d=|X_{d-2}|=|X_{(m-2)/2}|\;}
\]
以及
\[
\boxed{\;
\frac{|\mathcal Z_m|}{M_m}
=
\exp\!\left(-\frac{m}{2}\log\varphi+o(m)\right).\;}
\]
因此，六倍窗零块具有严格的双相尺度：其绝对规模已达到半尺度稳定型空间的大小，而其相对密度却以同一半尺度斜率衰减。
\end{theorem}

\begin{proof}
由推论 \ref{cor:fold-zero-half-index-multiple6}，
\[
|\mathcal Z_m|
\ge
F_{(m+2)/2}
=
F_d.
\]
又因
\[
|X_n|=F_{n+2},
\]
取 \(n=d-2=(m-2)/2\) 即得
\[
F_d=|X_{d-2}|=|X_{(m-2)/2}|.
\]

另一方面，定理 \ref{thm:fold-zero-window6-density-sharp-exponent} 已证明
\[
\lim_{m\to\infty,\ m+2\equiv 0\!\!\!\pmod 6}
\frac{1}{m}\log\frac{|\mathcal Z_m|}{F_{m+2}}
=
-\frac12\log\varphi.
\]
这正等价于
\[
\frac{|\mathcal Z_m|}{M_m}
=
\exp\!\left(-\frac{m}{2}\log\varphi+o(m)\right).
\]
证毕。
\end{proof}

\begin{theorem}[window-\texorpdfstring{$10$}{10} 三扇区的例外塔刚性闭合]\label{thm:derived-window10-exceptional-trisector-rigid-closure}
在 \(m=10\) 的 bin-fold 退化谱中，存在唯一三扇区骨架
\[
\bigl(|\mathcal S_{10,5}|,\ |\mathcal S_{10,8}|,\ |\mathcal S_{10,9}|\bigr)=(55,52,37),
\]
并满足刚性闭合关系
\[
52=\dim\mathfrak f_4,
\qquad
52+D_{13}=52+26=78=\dim\mathfrak e_6,
\]
\[
D_{16}=55=\dim\mathfrak e_7-\dim\mathfrak e_6,
\qquad
D_{18}+D_{13}=89+26=115=\dim\mathfrak e_8-\dim\mathfrak e_7,
\]
以及
\[
D_{18}=89=52+37.
\]
因此，\(52\)-扇区是例外核扇区，\(55\)-扇区是完备化尺度扇区，而 \(37\)-扇区被强制成为补余接口扇区。
\leanverified{paper\_derived\_window10\_exceptional\_trisector\_rigid\_closure}
\end{theorem}

\begin{proof}
由段落 \ref{par:gut-trisector-55-52-37}，
\(m=10\) 的退化谱在稳定类型层强制出三扇区骨架
\[
\bigl(|\mathcal S_{10,5}|,\ |\mathcal S_{10,8}|,\ |\mathcal S_{10,9}|\bigr)=(55,52,37).
\]
同一小节的结论 \ref{par:gut-exceptional-kernel-f4} 给出
\[
|\mathcal S_{10,8}|=52=\dim\mathfrak f_4,
\]
而结论 \ref{par:gut-e6-completion-52-26} 则给出
\[
52+D_{13}=52+26=78=\dim\mathfrak e_6.
\]
再由结论 \ref{par:gut-maxfiber-exceptional-tower}，
\[
D_{16}=55=\dim\mathfrak e_7-\dim\mathfrak e_6,
\qquad
D_{18}+D_{13}=89+26=115=\dim\mathfrak e_8-\dim\mathfrak e_7.
\]
最后，接口化结论 \ref{par:gut-d18-endogenous-split} 给出
\[
D_{18}=|\mathcal S_{10,8}|+|\mathcal S_{10,9}|=52+37=89.
\]
因此三块 \(55/52/37\) 在同一层上完成了例外塔的 rigid closure。证毕。
\end{proof}

\begin{theorem}[window-\texorpdfstring{$10$}{10} 三扇区在资源层上彼此正交]\label{thm:derived-window10-trisector-resource-orthogonality}
\leanverified{paper\_derived\_window10\_trisector\_resource\_orthogonality}
沿用定理 \ref{thm:derived-window10-exceptional-trisector-rigid-closure} 的记号。则三扇区分别支配三种互不等价的资源：
\begin{enumerate}
\normalfont
  \item \(55\)-扇区支配完备化增量，因为
  \[
  55=D_{16}=\dim\mathfrak e_7-\dim\mathfrak e_6.
  \]
  \item \(52\)-扇区支配局部奇偶荷容量：在 \(m=10\) 的奇偶荷交换化中，\(d=8\) 的 \(52\) 条纤维产生一个秩 \(52\) 的 \((\ZZ/2\ZZ)\)-符号子晶格。
  \item \(37\)-扇区支配规范群熵的凸权重：尽管它只占类型份额
  \[
  \frac{37}{144},
  \]
  但其对规范群熵
  \[
  H_{10}^{\mathrm{gauge}}
  =
  55\log_2(5!)+52\log_2(8!)+37\log_2(9!)
  \]
  的贡献比例为
  \[
  \frac{37\log_2(9!)}{H_{10}^{\mathrm{gauge}}}\approx 0.3676,
  \]
  显著高于其类型占比。
\end{enumerate}
因此，三扇区并不是同一资源的三段，而分别支配完备化模、奇偶荷容量与规范熵凸权重三条正交资源轴。
\end{theorem}

\begin{proof}
第一条正是定理 \ref{thm:derived-window10-exceptional-trisector-rigid-closure} 的一部分。

第二条来自注记 \ref{rem:fold-bin-gauge-entropy-scale}：在 \(m=10\) 的退化谱
\[
\{5:55,\,8:52,\,9:37\}
\]
下，\(d=8\) 的 \(52\) 条纤维在奇偶荷交换化中产生一个秩 \(52\) 的偶阶符号子晶格，刻画可观测 \((\ZZ/2\ZZ)\)-奇偶通道的最大局部容量。

第三条同样由注记 \ref{rem:fold-bin-gauge-entropy-scale} 给出。按 bits 记号，
\[
H_{10}^{\mathrm{gauge}}
=
\sum_{w\in X_{10}}\log_2\bigl(d_{10}^{\mathrm{bin}}(w)!\bigr)
=
55\log_2(5!)+52\log_2(8!)+37\log_2(9!),
\]
而 \(d=9\) 扇区的贡献恰为 \(37\log_2(9!)\)，故其相对权重为
\[
\frac{37\log_2(9!)}{H_{10}^{\mathrm{gauge}}}\approx 0.3676.
\]
与此同时，该扇区的类型占比仅为 \(37/144\approx 0.2569\)。由于 \(\log(d!)\) 对 \(d\) 严格凸，较高退化度扇区在熵权重下必然超占优。证毕。
\end{proof}

\begin{corollary}[奇偶荷完备与乘法 faithful 之间存在不可修补裂口]\label{cor:derived-parity-complete-not-multiplicative-faithful}
\leanverified{paper\_derived\_parity\_complete\_not\_multiplicative\_faithful}
设分辨率 \(m\) 满足
\[
d_{\min}(m)\ge 2.
\]
则最小完备奇偶荷签名秩满足
\[
N_m=|X_m|.
\]
特别地，在 \(m=6\) 时，
\[
|X_6|=21,
\qquad
\bigl(G_6^{\mathrm{bin}}\bigr)^{\mathrm{ab}}\cong (\ZZ/2\ZZ)^{21}.
\]
因此，任何试图把可观测异常压缩到少于 \(21\) 个独立比特的方案都必然破坏纤维独立置换结构。

然而，由定理 \ref{thm:derived-finite-anomaly-ledger-no-faithful-multiplicative-skeleton}，任何有限秩异常账本又都不能在严格同态语义下 faithful 承载乘法骨架。故
\[
\text{parity-complete}\not\Rightarrow \text{multiplicative-faithful},
\qquad
\text{finite anomaly ledger}\not\Rightarrow \text{arithmetic ledger}.
\]
\end{corollary}

\begin{proof}
由定理 \ref{thm:xi-foldbin-parity-charge-split-exact-minrank}，
\[
N_m=\#\{w\in X_m:\ d_m^{\mathrm{bin}}(w)\ge 2\}.
\]
若 \(d_{\min}(m)\ge 2\)，则每条纤维都满足 \(d_m^{\mathrm{bin}}(w)\ge 2\)，故
\[
N_m=|X_m|.
\]
在 \(m=6\) 时，由同一定理及 window-\(6\) 的直方图 \(2:8,3:4,4:9\)，得到
\[
|X_6|=21,
\qquad
\bigl(G_6^{\mathrm{bin}}\bigr)^{\mathrm{ab}}\cong (\ZZ/2\ZZ)^{21}.
\]
因此少于 \(21\) 个独立 \((\ZZ/2)\)-坐标不可能 complete 地恢复奇偶荷读出。

另一方面，定理
\ref{thm:derived-finite-anomaly-ledger-no-faithful-multiplicative-skeleton}
说明：任何有限秩异常账本都不可能在严格同态语义下 faithful 承载 \(\NN_{\ge1}^{\times}\) 的乘法结构。
两者合并即得所述裂口。证毕。
\end{proof}

\begin{conjecture}[faithful \texorpdfstring{$\pi$}{pi}-formula 的零块--例外核--乘法账本三寄存器原理]\label{conj:derived-pi-faithful-zero-exc-mult-three-register}
任一若要在这条支链上 faithful 地承载 projective \(\pi\)-通道，不应只依赖单一值层标量，而至少应同时外置三类彼此正交的寄存器
\[
\mathfrak R_{\mathrm{zero}},
\qquad
\mathfrak R_{\mathrm{exc}},
\qquad
\mathfrak R_{\mathrm{mult}}.
\]
其中：
\begin{enumerate}
\normalfont
  \item \(\mathfrak R_{\mathrm{zero}}\) 负责定理 \ref{thm:derived-window6-zero-block-dual-halfscale-law} 的六倍窗零块双相尺度；
  \item \(\mathfrak R_{\mathrm{exc}}\) 负责定理 \ref{thm:derived-window10-exceptional-trisector-rigid-closure} 与定理 \ref{thm:derived-window10-trisector-resource-orthogonality} 的例外三扇区分工，即 \(52\)-核、\(55\)-完备模与 \(37\)-接口余量；
  \item \(\mathfrak R_{\mathrm{mult}}\) 负责定理 \ref{thm:derived-finite-anomaly-ledger-no-faithful-multiplicative-skeleton} 与推论 \ref{cor:derived-parity-complete-not-multiplicative-faithful} 所禁止压缩的乘法 prime-ledger。
\end{enumerate}
该猜想表达的是：零块的半尺度二相性不能由例外三扇区资源替代；三扇区的正交分工又不能由奇偶荷签名替代；而有限奇偶荷签名更不能承载乘法 faithful。故任何试图把三者再压缩为单一数值公式的机制，至多是值层算法，而不可能是结构 faithful 的 \(\pi\)-completion 证书。
\end{conjecture}

\begin{remark}[统一链]\label{rem:derived-zero-exceptional-primeledger-unified-chain}
定理 \ref{thm:derived-finite-anomaly-ledger-no-faithful-multiplicative-skeleton}、
定理 \ref{thm:derived-window6-zero-block-dual-halfscale-law}、
定理 \ref{thm:derived-window10-exceptional-trisector-rigid-closure}、
定理 \ref{thm:derived-window10-trisector-resource-orthogonality}
与推论 \ref{cor:derived-parity-complete-not-multiplicative-faithful}
共同表明：在这条支链上，零块规模、例外塔补全与乘法账本并不处于同一压缩层。它们分别驻留在半尺度谱层、例外核接口层与不可有限同态化的 prime-ledger 层，因此任何 faithful completion 都必须显式保留这三层中的至少一个独立寄存器方向。
\end{remark}

\endinput
